\chapter{Beyond the Amateur Bands: Aviation, Maritime, and Military Radio}
\label{ch:beyond-amateur}

The radio spectrum an amateur studies for exams is one neighborhood in
a much larger city. Cheap, wideband software-defined radio hardware ---
a USB dongle costing little more than a cable, paired with free
decoding software --- has made it trivial for anyone with a computer to
\emph{listen} across aviation, maritime, and other services that once
required specialized, expensive receivers. This chapter is a
broadening excursion outside the amateur allocations: it will not be on
your license exam, but it uses exactly the same physics (modulation,
propagation, antennas) developed everywhere else in this book, and it
is one of the most popular on-ramps into the hobby for newcomers
arriving via SDR rather than via a first callsign.

\section{A Note on Legality Before You Start}

\begin{safetynote}[Receiving is not the same everywhere]
Whether it is legal to \emph{receive} (listen to) transmissions outside
the amateur bands --- and whether it is legal to \emph{disclose} the
content of what you hear --- varies enormously by country, and
sometimes by the specific service involved. Many countries permit
receive-only monitoring of aviation and maritime traffic freely; others
restrict it; several restrict listening to cellular telephone traffic
specifically even where other monitoring is permitted; a smaller number
restrict radio scanning broadly. \textbf{Transmitting} on any frequency
outside your license privileges is illegal essentially everywhere and
is not discussed in this chapter. Before monitoring any service
described here, confirm the specific rules in your own country. This
chapter describes what these signals \emph{are} and how they work,
purely as radio science and engineering education --- it is not
guidance on where such monitoring is or is not legally permitted.
\end{safetynote}

\section{The SDR Monitoring Ecosystem}

The same direct-sampling and quadrature-sampling architectures
described in Chapter~\ref{ch:receivers} underpin an entire hobby
ecosystem built around inexpensive wideband receivers (commonly
covering roughly 24\,MHz to 1.7\,GHz or wider on the least expensive
devices, extended further with upconverters or purpose-built HF-capable
units) feeding open-source decoding software. Because these receivers
present raw digitized RF to a host computer rather than a single fixed
demodulator, one piece of hardware can be retargeted from one service to
another purely by changing software --- the same receiver that tracks
aircraft one moment can decode marine safety broadcasts the next.
Antenna choice matters considerably: a simple dipole cut for one amateur
band performs poorly far outside that band, so serious multi-service
monitoring typically uses a wideband discone, log-periodic, or
service-specific antenna (an airband-optimized vertical, for instance)
rather than a repurposed amateur antenna.

\section{Aviation}

\subsection{VHF Air-Band Voice}

Civil aviation voice communication (air traffic control, pilot
position reports) uses AM voice --- notably \emph{not} FM ---
in the roughly 118--137\,MHz band, internationally, with 25\,kHz or (in
some regions) narrower 8.33\,kHz channel spacing. AM was retained for
aviation long after FM became dominant elsewhere in VHF land mobile use
for a specific safety reason: when two aircraft transmit AM signals on
the same channel simultaneously, both are typically still at least
partially intelligible (a heterodyne tone is heard, but speech often
remains readable); simultaneous FM transmissions on the same channel
instead tend to produce the FM capture effect (Chapter~\ref{ch:modulation}),
where only the stronger signal is heard and the weaker one is
suppressed entirely --- far more dangerous if the suppressed
transmission was a safety-critical call.

\subsection{ACARS and VDL}

\textbf{ACARS} (Aircraft Communications Addressing and Reporting
System) is a digital datalink that airliners use to exchange short,
structured messages with airline operations centers --- position
reports, weather requests, maintenance data, gate information --- largely
automatically, without a pilot manually keying a microphone. Classic
ACARS uses a relatively simple 2400-baud audio-frequency-shift-keyed
signal on VHF frequencies around 131\,MHz, decodable with widely
available free software directly from an SDR's audio output; a modern
successor, \textbf{VDL Mode~2}, uses a more spectrally efficient
D8PSK-modulated signal carrying substantially more traffic and is now
the dominant form on the primary international ACARS frequency.
Decoded ACARS traffic is largely operational/administrative rather than
safety-of-life air traffic control communication, but it offers a vivid,
approachable illustration of a real-world FSK/PSK digital data link
(Chapter~\ref{ch:modulation}) running continuously overhead.

\subsection{Mode S, ADS-B, and Aircraft Tracking}

Modern aircraft transponders reply to ground radar interrogation, and
increasingly broadcast unprompted, on \textbf{1090\,MHz} using a
pulse-position-modulated format. \textbf{ADS-B} (Automatic Dependent
Surveillance--Broadcast) extends this into a self-announced digital
message containing an aircraft's GPS-derived position, altitude,
velocity, and identity, transmitted roughly once per second without
requiring any ground interrogation. Because 1090\,MHz ADS-B messages are
unencrypted and specifically designed for wide reception, decoding them
with an inexpensive SDR dongle and free software is one of the most
popular and beginner-friendly SDR projects, commonly used to build
real-time aircraft tracking maps and to contribute data to
crowd-sourced flight-tracking networks.

\subsection{HF Aviation}

Over oceans and remote regions beyond VHF line-of-sight range,
aviation falls back to HF skywave propagation
(Chapter~\ref{ch:propagation}), using dedicated internationally
coordinated frequencies for long-range air traffic control voice and,
increasingly, the digital \textbf{HFDL} (HF Data Link) datalink ---
functionally an HF-band counterpart to VHF ACARS/VDL, using multiple
simplex ground-station frequencies with the aircraft selecting whichever
currently propagates best. \textbf{SELCAL} (selective calling) tones
allow a ground station to silently alert a specific aircraft's flight
deck on an otherwise-monitored HF voice channel without requiring
continuous crew listening.

\section{Maritime}

\subsection{VHF Marine Band and Digital Selective Calling}

Maritime VHF (approximately 156--174\,MHz, internationally standardized
into numbered channels) carries ship-to-ship and ship-to-shore voice
traffic, with specific channels reserved by international agreement for
distress, safety, and calling purposes. \textbf{Digital Selective
Calling (DSC)}, a burst digital signal on a dedicated calling channel,
allows a vessel to transmit an automated, GPS-tagged distress alert
--- identifying the vessel and its position --- at the press of a single
guarded button, then directs other stations to a working voice channel,
substantially speeding up distress response compared to voice alerting
alone.

\subsection{NAVTEX}

\textbf{NAVTEX} broadcasts maritime safety information --- navigational
warnings, weather forecasts and warnings, search-and-rescue
notices --- as narrow-band direct-printing (digital text) transmissions
on standardized frequencies (internationally, principally 518\,kHz),
automatically received and printed or displayed by dedicated shipboard
receivers without requiring an operator to be actively listening,
extending the same store-and-forward safety philosophy amateur emergency
nets rely on (Chapter~\ref{ch:emergency-comms}) into routine commercial
and leisure maritime operation.

\subsection{AIS: Automatic Identification System}

\textbf{AIS} is the maritime counterpart to aviation's ADS-B: vessels
above a certain size (and many smaller ones voluntarily) broadcast
position, course, speed, and identity using a GMSK-modulated digital
signal on two dedicated VHF channels (161.975 and 162.025\,MHz),
enabling both collision avoidance between nearby vessels and
shore-based or satellite-based global vessel tracking. Like ADS-B, AIS
is unencrypted by design and is a popular, approachable SDR decoding
project.

\subsection{The GMDSS Framework}

These maritime systems --- VHF DSC, NAVTEX, and additional HF/satellite
components not detailed here --- fit within the international
\textbf{Global Maritime Distress and Safety System (GMDSS)}, a
layered, largely automated framework designed so that a vessel in
distress has multiple independent, largely automatic paths to alert
rescue authorities even if a radio operator cannot personally send a
manual distress call --- a philosophy directly comparable to the
redundancy and priority-of-life-safety-traffic principles discussed for
amateur emergency communications in Chapter~\ref{ch:emergency-comms}.

\section{Military and Government Radio}

Military and government radio use spans the same physics as every
other service in this book, generally with three practical differences
from civil use: greater use of encryption and specialized waveforms
(which SDR monitoring cannot and should not attempt to defeat), greater
use of frequency-hopping and other anti-jam/anti-intercept techniques,
and, for what remains observable, a strong tradition of extremely
disciplined, standardized procedure.

\subsection{HF Global Communications}

Long-range military HF communication, like civil aviation and maritime
HF, relies on ionospheric skywave propagation
(Chapter~\ref{ch:propagation}). Some military HF networks (for example,
the United States Air Force's High Frequency Global Communications
System, historically monitored under callsigns beginning ``Skyking'' for
certain emergency-action traffic) use plain, highly standardized voice
procedure specifically so that the format itself, not secrecy of the
channel, provides authentication --- an interesting contrast to the
assumption that ``military'' automatically means ``encrypted.''

\subsection{Numbers Stations}

\textbf{Numbers stations} --- shortwave broadcasts reading strings of
spoken digits, letters, or coded phrases, often preceded by a
distinctive interval signal or melody --- have been observed on HF for
decades and are widely believed (though rarely officially acknowledged
by any government) to transmit one-time-pad-encrypted messages to field
intelligence operatives, since a correctly used one-time pad is
theoretically unbreakable regardless of how many people listen to the
broadcast. Whatever their true purpose, numbers stations are a
freely and legally receivable, genuinely strange corner of the HF
spectrum, and a favorite target for SDR hobbyists purely as a listening
curiosity.

\subsection{Trunked and Digital Land Mobile Radio}

Public safety and some military land-mobile users increasingly use
\textbf{trunked radio} systems (where a pool of frequencies is shared
dynamically among many talkgroups under computer control, rather than
each group having a permanently dedicated channel) built on digital
standards such as P25 or DMR (Chapter~\ref{ch:digital-modes} covers the
general digital voice/vocoder concepts these systems share with amateur
digital voice). Many such systems are receivable in the clear with
appropriate SDR software; increasingly many others are encrypted, which
SDR software can identify (the traffic is present and structured) without
being able to decode the actual content.

\section{Why This Matters Even If You Never Transmit Here}

\begin{keyidea}
Every signal in this chapter is built from the same building blocks
covered earlier in this book: a carrier, a modulation scheme
(Chapter~\ref{ch:modulation}), a propagation path
(Chapter~\ref{ch:propagation}), and, for the digital ones, an error-control
scheme (Chapter~\ref{ch:error-control}). Listening to AM aviation voice
next to FM amateur repeater traffic, or watching an ADS-B decoder plot
aircraft in real time next to an APRS map (Chapter~\ref{ch:digital-modes}),
is one of the fastest ways to make those chapters feel like real,
observable physics rather than abstract formulas --- which is exactly
why this material, though absent from any license exam, belongs in a
book about understanding radio deeply rather than just passing a test.
\end{keyidea}

\begin{quickreview}
\item Legality of receive-only monitoring, and of disclosing what is
received, varies by country and by service --- verify local rules before
monitoring; transmitting outside your license privileges is illegal
essentially everywhere.
\item Cheap wideband SDR hardware plus free software lets one receiver
retarget across aviation, maritime, and other services entirely in
software.
\item Aviation uses AM voice (deliberately, for graceful co-channel
degradation), ACARS/VDL digital datalink, and unencrypted 1090\,MHz
ADS-B position broadcasts; HF extends aviation communication beyond
VHF range with HFDL and SELCAL.
\item Maritime uses VHF voice with DSC digital distress alerting,
NAVTEX safety broadcasts, and AIS vessel tracking, all coordinated
under the GMDSS framework.
\item Military/government use spans the same physics with more
encryption, frequency-hopping, and standardized procedure; numbers
stations and observable-but-encrypted trunked systems are notable SDR
curiosities.
\end{quickreview}

\begin{practiceproblems}
\item Explain why civil aviation retained AM voice rather than
switching to FM, in terms of co-channel interference behavior.
\item Explain the functional difference between classic ACARS and
VDL~Mode~2, referencing the modulation concepts of
Chapter~\ref{ch:modulation}.
\item Explain why ADS-B and AIS are described as ``unencrypted by
design,'' and why that design choice makes sense given their safety
purpose.
\item Explain the role of DSC in maritime distress alerting, and how it
differs from a traditional spoken Mayday call.
\item Before monitoring any service in this chapter, what should you
personally verify first, and why does this book decline to state a
single universal answer?
\end{practiceproblems}
